\documentclass[11pt,a4paper]{article}
\usepackage{fontspec}
\setmainfont{DejaVu Sans}[Scale=0.90]
\usepackage{amsmath,amssymb}
\usepackage[margin=2.4cm]{geometry}
\usepackage{booktabs,array,longtable,tabularx}
\usepackage{enumitem}
\usepackage{titlesec}
\usepackage{parskip}
\titlelabel{\thetitle.\quad}
\titleformat{\section}{\normalfont\large\bfseries}{\thesection.}{0.6em}{}
\titleformat{\subsection}{\normalfont\normalsize\bfseries}{\thesubsection}{0.6em}{}
\titlespacing*{\section}{0pt}{16pt}{6pt}
\titlespacing*{\subsection}{0pt}{12pt}{4pt}
\setlist[itemize]{leftmargin=1.4em,itemsep=1pt,topsep=3pt}
\newcolumntype{L}[1]{>{\raggedright\arraybackslash}p{#1}}
\begin{document}
\begin{center}
{\Large\bfseries Fiche de lecture --- Controverse d'identifiabilité de 2023}\\[4pt]
{\small Krapez, \emph{J. Appl. Phys.} \textbf{134}, 056101 (2023) ; réponse de Mandelis, Kooshki \&
Melnikov, \emph{J. Appl. Phys.} \textbf{134}, 056102 (2023)}\\[2pt]
{\small Version 1 --- 12 septembre 2026}
\end{center}

\vspace{4pt}

\section{Identification}

\begin{center}
\begin{tabularx}{\textwidth}{@{}lX@{}}
\toprule
Article initial & A. Mandelis, S. Kooshki, A. Melnikov, \emph{Simultaneous density and thermal
conductivity depth profile reconstructions from noised thermal-wave amplitude and phase data using a
combined integral-equation and imperialist competitive algorithm method}, \emph{Journal of Applied
Physics} \textbf{133}, 055102 (2023), doi:10.1063/5.0129536 \\
\addlinespace
Commentaire & J.-C. Krapez (ONERA, DOTA), \emph{Journal of Applied Physics} \textbf{134}, 056101 (2023),
doi:10.1063/5.0151465, accès libre \\
\addlinespace
Réponse & A. Mandelis, S. Kooshki, A. Melnikov, \emph{Journal of Applied Physics} \textbf{134}, 056102
(2023), doi:10.1063/5.0155055 \\
\addlinespace
Nature & Échange critique en trois pièces : méthode d'inversion, objection d'identifiabilité, réponse
des auteurs \\
\bottomrule
\end{tabularx}
\end{center}

\section{Statut documentaire}

Le texte intégral du Commentaire est en accès libre et a servi de base à la présente fiche. Le texte de
la Réponse relève de l'accès sur abonnement et n'a pas pu être consulté lors de la rédaction : seules
les données bibliographiques en ont été vérifiées. Aucun résumé de son argumentation n'est donc proposé
ici. La section 8 se limite à situer la pièce dans l'échange et à énoncer les points sur lesquels une
réfutation devrait porter pour être recevable.

Les développements mathématiques des sections 5 et 6 sont des reconstructions établies dans les
conventions du projet. Ils fournissent la structure de l'objection et se substituent, en l'état, à toute
paraphrase des textes originaux.

\section{Objet de la controverse}

L'article initial présente une méthode d'inversion associant une formulation par équation intégrale du
problème aux limites de l'onde thermique et un algorithme d'optimisation globale de type compétition
impérialiste. La méthode est appliquée à un échantillon de métallurgie des poudres frittées présentant
une couche superficielle plus dense que le volume, et revendique la reconstruction \emph{simultanée} des
profils de masse volumique et de conductivité thermique en profondeur, la chaleur massique étant
supposée connue et constante.

Le Commentaire oppose que, pour la configuration thermophysique considérée, les profils de conductivité
thermique et de capacité thermique volumique ne peuvent pas être identifiés simultanément, parce qu'ils
sont pleinement corrélés. Le problème présente une non-identifiabilité inconditionnelle, d'où la
non-unicité de la solution. Le succès apparent des tests d'inversion n'y change rien : pour toute
solution apparemment optimale, on peut construire une infinité de couples de profils conjugués
également satisfaisants. L'identification simultanée des deux profils est donc tenue pour un problème
inverse insoluble.

Le Commentaire indique enfin que la transformation de Liouville éclaire l'influence des deux profils sur
la réponse en température, met en évidence le rôle propre de leur produit, c'est-à-dire de l'effusivité,
et révèle par là même la corrélation complète des deux profils. L'argument s'appuie sur des travaux
antérieurs de l'auteur consacrés à la modélisation du transfert thermique en milieu hétérogène gradué.

\section{Rappels physiques}

\subsection{Milieu gradué et grandeurs en jeu}

Le milieu est décrit par deux fonctions de la profondeur, la conductivité $\lambda(z)$ et la capacité
thermique volumique $\rho c(z)$, dont se déduisent
\[ a(z) = \frac{\lambda(z)}{\rho c(z)} , \qquad b(z) = \sqrt{\lambda(z)\,\rho c(z)} . \]
Lorsque la chaleur massique $c$ est supposée connue et constante, le profil de masse volumique se déduit
immédiatement de $\rho c$ par division. Une indétermination portant sur $\rho c$ se transporte donc
telle quelle sur la masse volumique.

\subsection{Configuration expérimentale}

L'expérience considérée consiste à appliquer un flux périodique à la surface d'un matériau opaque et à
analyser la réponse en température de cette même surface, par balayage en fréquence de l'amplitude et
de la phase. Le transfert est unidimensionnel, la loi de conduction est celle de Fourier, il n'y a pas
de sources internes, et l'observable est une fonctionnelle des seules propriétés du milieu.

\subsection{Équation à traiter}

Après transformée temporelle et condition initiale nulle,
\[ \frac{d}{dz}\!\left[\lambda(z)\,\frac{d\theta}{dz}\right] = p\,\rho c(z)\,\theta ,
\qquad \varphi = -\lambda(z)\,\frac{d\theta}{dz} . \]

\section{Rappels mathématiques}

\subsection{Transformation de Liouville}

Avec la coordonnée transformée et la variable de flux
\[ \xi(z) = \int_{0}^{z}\frac{du}{\sqrt{a(u)}} , \qquad \phi = -\,b(\xi)\,\theta' , \]
l'équation précédente devient
\[ \frac{d}{d\xi}\!\left[b(\xi)\,\frac{d\theta}{d\xi}\right] = p\,b(\xi)\,\theta . \]

La vérification est directe : $d/dz = a^{-1/2}\,d/d\xi$, d'où
$a^{-1/2}\,(b\,\theta_{\xi})_{\xi} = p\,\rho c\,\theta$, et l'identité $\rho c\,\sqrt{a} = b$ fournit le
résultat.

\subsection{Énoncé d'invariance}

Le point décisif tient en une phrase : \emph{dans la coordonnée transformée, une seule fonction
matériau subsiste, l'effusivité $b(\xi)$}. La réponse de surface, rapport de la température au flux en
$\xi=0$, est donc une fonctionnelle de $b(\cdot)$ seule. Les deux fonctions $\lambda$ et $\rho c$
n'interviennent que par cette combinaison et par le changement de variable qu'elles induisent.

La forme normale s'en déduit par la substitution habituelle :
\[ \frac{d^{2}\psi}{d\xi^{2}} - \bigl(V(\xi)+p\bigr)\psi = 0 ,
\qquad \psi = s\,\theta , \qquad s = b^{\pm1/2} , \qquad V = \frac{s''}{s} . \]

\subsection{Famille des profils conjugués}

Soit $\hat z(\xi)$ une fonction quelconque, croissante et dérivable, telle que $\hat z(0)=0$. On pose
\[ \tilde a(\hat z) = \left(\frac{d\hat z}{d\xi}\right)^{\!2} , \qquad
\tilde b(\hat z) = b(\xi) , \]
d'où, par les relations inverses,
\[ \tilde\lambda = b\,\frac{d\hat z}{d\xi} , \qquad
\widetilde{\rho c} = b\,\frac{d\xi}{d\hat z} , \qquad
\tilde\lambda\,\widetilde{\rho c} = b^{2} = \lambda\,\rho c . \]

Le milieu ainsi construit possède, par construction, la même effusivité exprimée dans la même
coordonnée transformée. Sa réponse de surface est donc rigoureusement identique à celle du milieu
initial, pour toute fréquence et en l'absence de bruit.

La famille est paramétrée par une fonction arbitraire : l'indétermination est de dimension
fonctionnelle et non de dimension finie. Le produit $\lambda\,\rho c$ est conservé, tandis que le
rapport $\lambda/\rho c$ est entièrement libre.

\subsection{Version infinitésimale}

En posant $\hat z_{\varepsilon} = z + \varepsilon\,w$, la variation au premier ordre s'écrit, à
coordonnée transformée fixée,
\[ \frac{\delta\lambda}{\lambda} = \frac{w'}{z'} , \qquad
\frac{\delta(\rho c)}{\rho c} = -\frac{w'}{z'} , \qquad
\frac{\delta\lambda}{\lambda} + \frac{\delta(\rho c)}{\rho c} = 0 . \]
Le noyau de l'application tangente est donc non réduit à zéro : toute perturbation relative opposée des
deux profils, accompagnée du déplacement de profondeur correspondant, laisse l'observable inchangée.
La matrice d'information associée à une discrétisation quelconque des deux profils est singulière, quel
que soit le nombre de fréquences.

\subsection{Réduction au cas homogène}

Le choix particulier $\hat z = \mu z$ donne $d\hat z/d\xi = \mu\sqrt{a}$, d'où
\[ \tilde a = \mu^{2}a , \qquad \tilde\lambda = \mu\lambda , \qquad
\widetilde{\rho c} = \mu^{-1}\rho c , \qquad e \mapsto \mu\,e , \]
c'est-à-dire exactement le groupe à un paramètre de la fiche relative à la note de 2017. Le résultat de
2017 est donc le cas homogène du résultat de 2023 ; la relation d'Euler entre sensibilités y était la
forme infinitésimale de la même invariance.

\begin{center}
\begin{tabularx}{\textwidth}{@{}lXX@{}}
\toprule
& \textbf{Note de 2017} & \textbf{Commentaire de 2023} \\
\midrule
Objet & Couche homogène sur substrat & Milieu gradué semi-infini \\
\addlinespace
Inconnues & Trois scalaires : $L$, $a$, $\lambda$ & Deux fonctions : $\lambda(z)$, $\rho c(z)$ \\
\addlinespace
Quantités identifiables & $\xi_{1}=L/\sqrt{a}$ et le rapport d'effusivité & L'effusivité exprimée dans
la coordonnée transformée \\
\addlinespace
Dimension de l'indétermination & Un paramètre & Une fonction \\
\addlinespace
Forme de l'argument & Dépendance linéaire des vecteurs de sensibilité & Noyau non trivial de
l'application tangente \\
\bottomrule
\end{tabularx}
\end{center}

\subsection{Non-identifiabilité inconditionnelle}

La qualification d'inconditionnelle signifie que la dégénérescence ne dépend ni du jeu de fréquences,
ni du niveau de bruit, ni de la méthode d'optimisation. Elle se distingue du mauvais conditionnement,
qui affecte une application injective mais instable. Les deux défauts se cumulent ici : à
l'indétermination structurelle décrite ci-dessus s'ajoute la perte de sensibilité aux variations
profondes, elle-même dépendante de la bande de mesure.

Une conséquence pratique mérite d'être isolée. Un algorithme d'optimisation globale converge vers un
point de l'orbite d'invariance, et non vers le profil réel. La qualité de l'ajustement n'est donc pas un
indice d'exactitude : elle est identique en tout point de l'orbite.

\section{Structure de l'objection}

\begin{center}
\begin{tabularx}{\textwidth}{@{}lX@{}}
\toprule
Prémisse 1 & L'observable est la réponse en température de la face avant, sous excitation de flux en
face avant, en transfert unidimensionnel sans sources internes \\
\addlinespace
Prémisse 2 & Après transformation de Liouville, cette observable ne dépend que de l'effusivité exprimée
dans la coordonnée transformée \\
\addlinespace
Conclusion 1 & Deux profils de conductivité et de capacité volumique de même produit, moyennant le
reparamétrage de profondeur associé, sont indiscernables \\
\addlinespace
Conclusion 2 & L'identification simultanée est impossible ; la chaleur massique étant supposée connue,
l'indétermination se transporte sur la masse volumique \\
\addlinespace
Conclusion 3 & Le succès numérique de l'inversion ne constitue pas une réfutation, puisque tout point de
l'orbite produit le même résidu \\
\bottomrule
\end{tabularx}
\end{center}

\section{Levées possibles de l'indétermination}

L'invariance repose sur la liberté de reparamétrer la profondeur. Toute information imposant une échelle
de longueur indépendante la restreint.

\begin{center}
\begin{tabularx}{\textwidth}{@{}L{4.6cm}XX@{}}
\toprule
\textbf{Information ajoutée} & \textbf{Effet sur l'orbite} & \textbf{Portée} \\
\midrule
Un des deux profils connu & Réduction à un point & Cas usuel : capacité volumique supposée scalaire
connue \\
\addlinespace
Repère géométrique interne connu & Restriction par condition aux limites en $\hat z$ & L'indétermination
demeure fonctionnelle, avec extrémité fixée \\
\addlinespace
Diffusion latérale, configuration bidimensionnelle & Introduction d'une longueur transverse & Brise
l'invariance d'échelle en profondeur \\
\addlinespace
Sources internes dans la couche étudiée & Le champ source fixe la coordonnée physique & Brise
l'invariance \\
\addlinespace
Mesure en face arrière d'un échantillon d'épaisseur connue & Fixe l'extrémité de l'orbite & Réduit sans
supprimer l'indétermination \\
\bottomrule
\end{tabularx}
\end{center}

Aucune de ces voies ne relève d'un raffinement de l'algorithme d'inversion : toutes consistent à
modifier le dispositif ou à ajouter une mesure de nature différente.

\section{Position de la Réponse}

La Réponse constitue la troisième pièce de l'échange, publiée immédiatement à la suite du Commentaire
dans le même numéro. Son texte n'a pas été consulté ; aucune de ses affirmations n'est donc rapportée
ici.

Pour mémoire, une réfutation recevable de l'objection devrait porter sur l'une au moins des prémisses de
la section 6, à savoir :

\begin{itemize}
\item contester la prémisse 1, en montrant que le dispositif réel comporte une information non prise en
compte dans le modèle unidimensionnel à excitation et détection en face avant ;
\item contester la prémisse 2, en exhibant une dépendance de l'observable à une grandeur autre que
l'effusivité exprimée dans la coordonnée transformée ;
\item contester la conclusion 1, en exhibant une contrainte de régularité ou de positivité éliminant
tous les profils conjugués sauf un.
\end{itemize}

Cette section reste à compléter dès l'obtention du texte intégral de la Réponse.

\section{Rattachement au présent travail}

\subsection{Continuité avec les conventions du projet}

La coordonnée $\xi$, la forme normale et le potentiel $V=s''/s$ employés dans le Commentaire sont ceux
que fixe le fichier de conventions. L'objection se formule donc sans conversion préalable.

\subsection{Apport propre}

Le Commentaire fournit un cas où la non-unicité structurelle est explicite et constructive : l'orbite
d'invariance est exhibée, non seulement établie. Il documente, en outre, une confusion récurrente dans
la littérature de métrologie thermique entre la convergence d'un algorithme et l'unicité de la solution.

\subsection{Divergence de cadre}

Comme la note de 2017, l'analyse se place sous la loi de Fourier. Sous une loi à temps de relaxation
fini, le paramètre spectral devient $p+\tau p^{2}$ et la transformation de Liouville doit être reprise :
la question de savoir si l'invariance décrite en 5.3 subsiste, se déforme ou disparaît n'est pas traitée
par les textes de l'échange et demeure ouverte dans le présent travail.

\section{Références}

\begin{enumerate}[leftmargin=2.2em]
\item A. Mandelis, S. Kooshki, A. Melnikov, \emph{Journal of Applied Physics} \textbf{133}, 055102
(2023), doi:10.1063/5.0129536.
\item J.-C. Krapez, \emph{Comment on ``Simultaneous density and thermal conductivity depth profile
reconstructions from noised thermal-wave amplitude and phase data using a combined integral-equation and
imperialist competitive algorithm method'' [J. Appl. Phys. 133, 055102 (2023)]}, \emph{Journal of Applied
Physics} \textbf{134}, 056101 (2023), doi:10.1063/5.0151465.
\item A. Mandelis, S. Kooshki, A. Melnikov, \emph{Response to ``Comment on \ldots'' [J. Appl. Phys. 134,
056101 (2023)]}, \emph{Journal of Applied Physics} \textbf{134}, 056102 (2023), doi:10.1063/5.0155055.
\item J.-C. Krapez, \emph{International Journal of Thermal Sciences} \textbf{136}, 182--199 (2018).
\item D. Maillet, S. André, J.-C. Batsale, A. Degiovanni, C. Moyne, \emph{Thermal Quadrupoles: Solving
the Heat Equation through Integral Transforms}, Wiley, 2000.
\end{enumerate}
\end{document}
